\chapter{引言}\label{ux5f15ux8a00}

\section{写在前面：致新人类}\label{ux5199ux5728ux524dux9762ux81f4ux65b0ux4ebaux7c7b}

\noindent 亲爱的新人类：

当你们看到这封信时，我或许已经在轰炸中散为灰烬，或是掩埋在无人知晓的废墟下，或是在中国西南高原的某个深山中度过平静的一生，然后在某棵歪脖子树下长眠。

这是一个永远背负着愧疚的留守者。此刻坐在即将关闭的地面控制中心。我决定写下这封信，但不知道从何落笔。

我今年六十八岁，是``火种计划''的总工程师。在送走百万立方球体前，我在杭州生活了五十年。我的妻子在三年前的大华东战争中去世。选择留在地球，是我一生中最痛苦的决定。

我们分别的那天，地球上还有二十亿人在挣扎求生。我记得在发射基地，你们的两千一百三十九分之一，透过舷窗望着我，他的眼神我至今难忘。我们都是人类，都有生存的权利，但我们这些所谓的``决策者''却要决定谁走谁留，这个画面可能会伴随着我直到我不算太长的余生的尽头。

旧文明最后的日子让我明白了很多事。我们曾经拥有那么多：完善的医疗体系，便捷的交通网络，丰富的文化传承。但我们同样也犯下了太多错误。我亲眼见过妇幼医院因为轰炸而断电，见过温和的邻居为了一箱纯净水而互相砍杀，见过世界上最顶尖的科学家们为国家的利益而集体销毁数据。\\


我想告诉你们一些很重要的事：

\begin{itemize}
\item
  要建立透明的预警机制。不要因为任何理由隐瞒真相，再残酷的事实也比美丽的谎言更有价值。
\item
  要珍惜你们的共同体。分歧与立场永远存在，但不要让分歧演变成仇恨。我们那个时代，人们为了国家立场、宗教信仰、民族矛盾争论不休，最后却忘记了我们都是人类。
\item
  要对权力保持警惕。任何绝对的权力都会导向绝对的奴役，需要让权力握在每个人手中，需要建立有效的制衡机制。
\item
  要重视基础科学。我们太注重短期利益，太多优秀的头脑去研究如何制造泡沫，蛀空实体经济，而不是解决人类面临的真正问题。\\
\end{itemize}

我想起今早最后一次去曾经被称作超市的地方，货架上空空如也；想起邻居家的钢琴声，已经很久没有响起；想起二十年前的那个春天，我和妻子去荷兰开会，那里有全世界闻名的郁金香海。那些郁金香现在应该或许还在开放，只是再也不属于尼德兰共和国了。

孩子们，你们生活在一个新的起点。请记住我们的教训，但不要被我们的罪孽压垮。要建立公平的制度，但要包容人性的弱点；要追求进步，但不要牺牲普通人的幸福；要发展科技，但不要忘记道德的约束。

我们这些留下的人，每个人都带着无法愈合的伤口。有的人在深夜哭泣，有的人一直沉默，有的人反复做着关于过去的噩梦。我们不是英雄，只是一群在绝境中选择了承担最后责任的普通人。

愿你们能建立更好的文明，愿你们能原谅我们的过错，愿你们永远不会再被迫面临我们这样的抉择。

再会。\\


又及：请保管好我送去的种子，它们是我家乡的品种，很古老，叶子呈扇形，很好看。

\rightline{------【前文明】``火种计划''发起者 Aleph\_null}

\newpage

\section{《日轮之冕序》}\label{ux65e5ux8f6eux4e4bux5195ux5e8f}

\begin{center}
鸿蒙初辟\,文明肇端

求真问道\,薪火相传

舟车万里\,文教四藩

寰宇同风\,盛世如磐

衅生末季\,祸起萧墙

天灾频仍\,人欲横张

利争不让\,义战沦亡

干戈四动\,血染玄黄

兄弟阋墙\,同室操戈

金城尽毁\,典籍沉疴

饥馑遍野\,疫病成魔

礼乐经纶\,几近倾磨

泣血以悼\,往圣先贤

断简残篇\,犹记华年

千秋功业\,尽化寒烟

万家灯火\,一炬成眠

今立新邦\,惕厉自省

前事不忘\,后事之师

克明俊德\,以制人欲

立宪均平\,以协众心

苍天为鉴\,后土为凭

文脉永续\,英华不朽

谨以斯文\,告慰先圣

继绝开新\,万世长存
\end{center}

\newpage

\section{新地球与文明的愿景}\label{ux65b0ux5730ux7403ux4e0eux6587ux660eux7684ux613fux666f}

\subsection{序言}\label{ux5e8fux8a00}

我全体文明之公民，

同兹决心，

免后世再遭旧世界崩解之浩劫，其苦痛已为人类共同之记忆，

重申对生存、尊严、平等与世代延续之基本信念，

创造稳定环境，以维持公义，恪守由共同命运而生之义务，

运用集体之力量，以促进文明之进步与全体福祉之提升。

爱由我全体公民之意志，

汇集于日轮之冕下，

确立此根本愿景，

作为我文明存续与发展之最高纲领。

\subsection{根本宗旨}\label{ux7b2cux4e00ux7ae0-ux6839ux672cux5b97ux65e8}

本文明之存续与发展，遵循以下最高宗旨：

\textbf{第一条}\,保障人类文明之整体安全与永续生存。采取一切必要之规划与行动，消除内外部对共同体存续之威胁，并以理性与团结之原则，应对任何可能破坏集体福祉之挑战。

\textbf{第二条}\,巩固并发展基于完全平等与共同奉献之公民关系。采取一切适当方式，增强公民间之纽带，培育牢不可破之共同体认同。

\textbf{第三条}\,借助统筹协作，解决关乎文明发展之经济、科技、环境与文化等根本性问题，并无差别地促进每一位公民之全面发展与其基本权利之实现。

\textbf{第四条}\,构成一个意志统一、行动协调之整体，以达致上述所有共同目标。

\subsection{核心原则}\label{ux7b2cux4e8cux7ae0-ux6838ux5fc3ux539fux5219}

为实现上述宗旨，本文明及其全体公民均应以下列原则为行事准则：

\textbf{第五条}\,本文明建立在所有公民权利与义务根本平等之基础上。

\textbf{第六条}\,全体公民应秉持最大之诚意，履行本愿景及据此衍生之一切法律所赋予之义务。

\textbf{第七条}\,文明内部之任何分歧，应以和平与协商之方式解决，不得危及共同之安全、团结与生存根基。

\textbf{第八条}\,全体公民在一切行动中，不得有任何违背本文明根本宗旨之行为，亦不得以任何形式损害共同体之整体利益。

\textbf{第九条}\,凡为维护文明整体安全与存续而采取之必要行动，全体公民应予遵从并为实现该等行动提供一切所需之支持。

\textbf{第十条}\,本愿景之精神与原则，应为全体公民永恒恪守之信条，并作为教育之核心，代代相传。

\subsection{终章}\label{ux7ec8ux7ae0}

\textbf{第十一条}\,此愿景所载之崇高理想与庄严原则，乃是我文明精神之不灭灯塔，必须广布于每一公民之心间，贯彻于文明之每一项事业。

\textbf{第十二条}\,我们，日轮之冕的公民，于此郑重宣誓：愿以我们全部的生命、智慧与勇气，共同捍卫、实践并传承此愿景。

为此誓言，

立此存照。
